
\lussec 5. 格点正规化

本节的出发点是 Wilson [\cite{Wilson}] 所提议的
把夸克场放到格点上的一种格点 QCD 表述。
对格点狄拉克算符或规范作用量不作特殊假设，
但二者应当是局域的，
并保持标准 Wilson 理论所保持的全部对称性。
为记号方便，取格距 $a=1$。
除非另有说明，格点在所有方向上被视为无限延伸。


\lussubsec 5.1 流方程

在格点上，随时间演化的规范场 $V(t,x,\mu)$ 由
\equation{
  \left.V\right|_{t=0}=U,
  \enum
  \nexteq{2.5ex}
  \partial_tVV^{-1}=-g_0^2\partial\Sw,
  \enum
}
定义，其中 $U(x,\mu)$ 是基本格点规范场，
$(\partial\Sw)(V,x,\mu)$ 是 Wilson 方格作用量的梯度
（进一步的解释见附录 A 及文献~[\cite{WilsonFlow}]）。

对于夸克场，边界条件与流方程的形式与连续理论相同，
只是式 (2.4) 中的算符 $\Delta$
换成格点拉普拉斯算子
\equation{
  \Delta=\nabstar{\mu}\nab{\mu},
  \enum
}
这里 $\nab{\mu}$ 与 $\nabstar{\mu}$
是在随时间演化规范场存在下的规范协变前向与后向差分算符。

如文献~[\cite{WilsonFlow}] 中已指出的，
流方程 (5.2) 解的整体存在性、唯一性与光滑性是被严格保证的。
由于 $\Delta$ 是有界算符，且它随流时间光滑变化，
可以证明同样的结论对夸克流方程也成立
[\cite{ArnoldI}]。
因此非零流时间的局域场——例如前面各节所考虑的那些场——
在格点上是完全良定义的。


\lussubsec 5.2 $4+1$ 维格点理论

与连续理论类似，
随时间演化局域场的格点 QCD 关联函数
可以从一个 $4+1$ 维格点场论得到。
为了能够完全控制这一构造，
先对该场论就流方程的一个离散化版本建立之，
其中流时间被限制在一个有限集合
\equation{
  0,\eps,2\eps,3\eps,\ldots,T
  \enum
}
上，各值由时间步长 $\eps>0$ 分隔。

离散演化方程
\equation{
  V(t+\eps,x,\mu)V(t,x,\mu)^{-1}=\rme^{-\eps g_0^2(\partial\Sw)(V,x,\mu)}
  \enum
}
对规范场的积分，相当于用欧拉方法
（或者等价地，通过施加一列 ``stout 涂抹'' 步骤 [\cite{Stout}]）
对连续方程 (5.2) 作近似积分。
对夸克场而言，
离散化的流方程形式与连续情形相同，
只是夸克场的时间导数换成前向差分
\equation{
  \partial_t\chi(t,x)={1\over\eps}\{\chi(t+\eps,x)-\chi(t,x)\}.
  \enum
}
注意这些离散方程可以逐步求解：
从 $t=0$ 的基本场出发，
递推地从时间 $t$ 推进到 $t+\eps$。
特别地，解由场的初值唯一确定。

在 $4+1$ 维格点理论中，
泛函积分中积分的所有场只在时刻 (5.4) 上定义。
由于此处不需要规范固定，
没有鬼场，作用量中也没有规范固定项，
但边界条件 (2.6)、(5.1) 照旧施加。
拉格朗日乘子场在时刻 $T$ 的分量与其他场退耦，
因此可置为零（或从一开始就略去）。

体作用量的一种可能选取是
\equation{
  \SGfl=-2\sum_{t=0}^{T-\eps}\sum_{x,\mu}
  \tr\bigl\{L(t,x,\mu)
  \noenum
  \nexteq{-0.5ex}
  {\phantom{\SGfl=-2\sum_{t=0}^{T-\eps}\sum_{x,\mu}}}
  \times
  \Psu\bigl(
  V(t+\eps,x,\mu)V(t,x,\mu)^{-1}-\rme^{-\eps g_0^2(\partial\Sw)(V,x,\mu)}
  \bigr)\bigr\},
  \enum
  \nexteq{2.5ex}
  \SFfl=\eps\sum_{t=0}^{T-\eps}\sum_x\bigl\{
  \lambdabar(t,x)\left(\partial_t-\Delta\right)\chi(t,x)
  +\chibar(t,x)\left(\lvec{\partial}_t-\lvec{\Delta}\right)\lambda(t,x)
  \bigr\},
  \enum
}
其中映射
\equation{
  M\to\Psu(M)={1\over2}(M-M^{\dagger})
  -{1\over2N}\kern1pt\tr(M-M^{\dagger})
  \enum
}
把任意复 $N\times N$ 矩阵映为 $\SUn$ 李代数的一个元素。
除这些作用量及格点 QCD 作用量外，
测度项
\equation{
  \Sms=\sum_{t=0}^{T-\eps}\sum_{x,\mu}
  {\cal L}_{\rm ms}(V(t+\eps,x,\mu)V(t,x,\mu)^{-1}),
  \enum
  \nexteq{3.0ex}
  {\cal L}_{\rm ms}(W)=
  \cases{-\ln\det(J(W)) & if $W\in\Dsu$,\cr
         \noalign{\vskip1ex}
         \infty         & otherwise,\cr}
  \enum
}
也必须包含在该格点理论的 total 作用量之中
（映射 (5.9) 的雅可比矩阵 $J(W)^{ab}$
及单位元的邻域 $\Dsu\subset\SUn$ 的定义见附录 B）。

给定了场及其作用量之后，
$4+1$ 维格点理论的泛函积分按标准方式定义。
注意测度项 (5.10)、(5.11)
实际上把泛函积分中的链变量限制在使
\equation{
  V(t+\eps,x,\mu)V(t,x,\mu)^{-1}\in\Dsu
  \quad\hbox{for all $t,x,\mu$}.
  \enum
}
成立的区域内。
由于作用量 $\SGfl$ 是纯虚的，
对拉格朗日乘子场 $L(t,x,\mu)$ 的积分不是绝对收敛的，
但应当理解，
由此可能产生的积分不确定性
通过在作用量中加入一个正比于拉格朗日乘子场平方的无穷小项来消除。


\lussubsec 5.3 四维理论的恢复

这样定义的、由场 $V,\chi,\chibar$ 构成的局域场的关联函数，
可以在一定程度上通过先对拉格朗日乘子场积分而求值。
由于 total 作用量对这些场是线性的，
积分给出一个乘积
\equation{
  \prod_{t=0}^{T-\eps}\prod_{x,\mu}
  \delta\bigl\{\Psu\bigl(
  V(t+\eps,x,\mu)V(t,x,\mu)^{-1}\bigr)-
  \Psu\bigl(\rme^{-\eps g_0^2(\partial\Sw)(V,x,\mu)}\bigr)\bigr\}
  \noenum
  \nexteq{1.0ex}
  \times\prod_{t=0}^{T-\eps}\prod_{x}
  \delta\bigl\{\partial_t\chi(t,x)-\Delta\chi(t,x)\bigr\}
  \delta\bigl\{\chibar(t,x)\lvec{\partial}_t-
  \chibar(t,x)\lvec{\Delta}\bigr\}
  \enum
}
即狄拉克 $\delta$ 函数的乘积。
当 $V,\chi,\chibar$ 满足离散流方程时，
这些 $\delta$ 函数的自变量为零；
而且对足够小的步长 $\eps$，
事实上没有别的办法能使这些自变量为零%
\kern1pt\sfn{$\dagger$}{%
若 $\exp\{\eps g_0^2(\partial\Sw)(V,x,\mu)\}\in\Dsu$
在格点的所有链上及所有流时间处成立，
则唯一性得到保证。
一些严格的估计表明：
若 $\eps\leq(24\sqrt{2}N)^{-1}$，则该条件被满足。}。
对规范场而言，
此处需要注意的要点是，
$\delta$ 函数中的映射 $\Psu$
作用在单位元小邻域内的 $\SUn$ 矩阵上。
附录 B 所引映射的性质于是意味着：
$\delta$ 函数所强加的约束等价于
把流方程 (5.5) 强加于规范场。

现在可以利用这些 $\delta$ 函数
对所有流时间 $t>0$ 处的场变量进行积分。
从时刻 $t=T$ 的场出发并向更小时间推进，
$\delta$ 函数允许递归地消去全部场，
直到剩下的积分变量只有基本场 $U,\psi,\psibar$。
在这一过程中，
规范场变量的 $\delta$ 函数积分产生映射 $\Psu$
的雅可比因子的乘积，
但这些因子被作用量中的测度项抵消
（见附录 B）。

计算于是表明：
不涉及拉格朗日乘子场的局域场的 $4+1$ 维关联函数，
与同一组场在四维中通过离散流方程和 QCD 泛函积分
定义的关联函数精确相等。


\lussubsec 5.4 连续时间极限

至此，关联函数中的随时间演化场仍是
由流方程的欧拉积分得到的那些场。
由于相关的积分误差为 $\eps$ 量级，
当 $\eps\to0$（且平凡地 $T\to\infty$）时
便恢复连续时间的关联函数。
此外，测度项
\equation{
  \Sms=\sum_{t=0}^{T-\eps}\sum_{x,\mu}
  \ln\det J\bigl(\rme^{-\eps g_0^2(\partial\Sw)(V,x,\mu)}\bigr)
  =\rmO(\eps)
  \enum
}
在此极限下为零。

下文中通常考虑格点理论的连续时间表述，
但 $4+1$ 维泛函积分总是理解为
通过离散时间泛函积分的极限 $\eps\to0$ 来定义。
